\section{Motivation}
\label{sec:motivation}
\label{sec:intro}

Suppose a server process embeds a SQL engine and lets its users format reports
with \sql{lpad}, which pads a string to a requested width. A user types a width
of 9{,}223{,}372{,}036{,}854{,}775{,}807. The server should answer with an error
message. If the engine instead ends the process, every other user's session ends
with it. A database engine is trusted to be correct without being checked, so a
defect of this kind is caught by the engine's own tests or not at all.

Database engines are among the most heavily tested programs in practice.
Differential testing runs the same query on two engines and reports any
disagreement. Metamorphic oracles compare a query with an equivalent rewrite of
itself~\cite{rigger2020tlp,rigger2020norec}. Grammar-based generators produce
valid SQL at scale~\cite{zhong2020squirrel,liang2026sqlbull}. All of these tools
stop at two boundaries: the generator can only emit the functions and values it
was written to emit, and a differential oracle can only judge a feature that its
reference engine also implements. We focus on the defects that sit past those
two boundaries. Such defects are silent in a specific way: a fuzzer can run for
a year without reaching them, and the time it has already run says nothing about
how close it is.

\begin{figure*}[t]
\begin{lstlisting}[escapechar=@]
-- (a) Regression test added by the fix for DuckDB issue #27846
select LPAD('x', 10, '0'), LPAD('x', 4, '@\textcolor{black}{\ding{72}}@'), LPAD('x', 3, '@\textcolor{black}{\ding{72}}@ab'), RPAD('x', 3, '@\textcolor{black}{\ding{72}}@ab')
-- (b) The same shape with every literal replaced by a typed hole
select LPAD({s4}, {n0}, {s5}), LPAD({s6}, {n1}, {s7}), LPAD({s8}, {n2}, {s9}), RPAD({s10}, {n3}, {s11})
-- (c) Filled from a boundary pool at seed 0, then run on Turso b588fbdc6
select LPAD('', -9223372036854775808, 'abc'), LPAD('_', -128, '%'), LPAD('0', 9223372036854775807, 'abc'), RPAD('%', -128, 'ABC')
-- (d) Reduced to one statement
select LPAD('a', 9223372036854775807, 'b');
@\textit{memory allocation of 9223372036854775806 bytes failed}@      @$\leftarrow$ process aborts, exit status 134@
\end{lstlisting}
\caption{The motivating example. A regression test from another engine~(a),
templated~(b) and refilled with boundary values~(c), aborts Turso from one
statement~(d). The marker \ding{72} stands for a four-byte emoji in the original
test. Turso's own fuzzers cannot produce~(d): their 47-function vocabulary has
no \sql{lpad}, and their integers stop at 1{,}000{,}000.}
\label{fig:motivating}
\end{figure*}

Figure~\ref{fig:motivating} shows a real case. Turso~\cite{turso} is a
SQLite-compatible engine written in Rust. It runs its own differential fuzzer in
continuous integration: 41{,}205 lines of Rust that generate SQL and compare
Turso's answers with SQLite's. Yet the statement in Figure~\ref{fig:motivating}d
ended the process with exit status 134 (SIGABRT) in 2.28 seconds, where it should
have returned the engine's own error \emph{string or blob too big}. Turso's
fuzzer could never have produced it. Its vocabulary of 47 scalar functions does
not contain \sql{lpad}, and it draws integer literals only from
$[-1{,}000{,}000,\ 1{,}000{,}000]$~\cite{tursofuzzer}. Even if it had produced
the statement, SQLite has no \sql{lpad} to compare the answer against. None of
Turso's four libFuzzer targets or its differential fuzzer calls \sql{lpad},
\sql{rpad} or \sql{repeat}.

The closest prior work either chooses where to look from the grammar or reuses
known defects through static analysis. SQLBull starts generation from grammar
rules that a tester marks as feature-rich~\cite{liang2026sqlbull}.
Recurring-vulnerability detectors turn a known patch into a rule and search for
code that breaks it~\cite{cao2025vulpa,cao2025antman}. However, the first needs
someone to pick the starting rules, and the second needs the defect to recur in
shared or copied code. Turso and DuckDB share no code.

We present \tool, which finds defects in this region in four steps:
\begin{enumerate}
  \item It reads the incumbent fuzzer's source and lists the functions and
        values that fuzzer cannot generate (Section~\ref{sec:step-blind}).
  \item It mines closed, reproduced bug reports from a different engine and keeps
        each report's query as a template whose literals are typed holes
        (Section~\ref{sec:step-mine}).
  \item It fills those holes, and the output of a typed SQL generator, with
        boundary values, and runs every statement against the engine under test
        (Section~\ref{sec:step-run}).
  \item It flags any statement that ends the process, then reduces, re-checks and
        reports it (Section~\ref{sec:step-triage}).
\end{enumerate}

We ran \tool against Turso. It found three previously unknown defects in Turso's
string functions: \sql{lpad} and \sql{rpad} allocate their result before checking
its length, and \sql{repeat} also overflows its size computation. Our patch for
them moved the length check in front of the allocation, and that exposed a
fourth defect, which \tool did not find: the check counts characters against a
limit in bytes. Our two patches, which fix all four, were merged
upstream~\cite{turso11225,turso11259}. At an equal budget of 20{,}000 cases, the
mined reports and the typed generator each reached both root causes of the
three defects. The lesson we draw is a practical one: before spending compute on
a fuzzer, spend an hour reading the fuzzer that is already there.

\paragraph{Research contributions.}
\begin{itemize}
  \item We define the \emph{blind spot} of an incumbent fuzzer and show that it
        can be measured from that fuzzer's source (Section~\ref{sec:problem}).
  \item We present \tool, which fills the blind spot with the shapes of another
        engine's bug reports and with typed generation, and judges results with an
        oracle that needs no reference engine (Section~\ref{sec:technique}).
  \item We report three previously unknown defects in Turso found by \tool, each
        regenerated from a recorded seed, and a fourth that our fix exposed. All
        four are fixed upstream by our merged patches (Section~\ref{sec:rq1}).
  \item We measure what each input source contributes at an equal budget, and
        release the tool, the corpus and all data
        (Sections~\ref{sec:rq2} and~\ref{sec:data}).
\end{itemize}

The rest of this section gives the background the paper needs
(Section~\ref{sec:background}), states the problem precisely
(Section~\ref{sec:problem}), and explains why the obvious alternatives do not
solve it.

\subsection{Background}
\label{sec:background}

\paragraph{Turso.} Turso~\cite{turso} reimplements SQLite in Rust and aims to
accept the same SQL. Its shell, \texttt{tursodb}, reads statements on standard
input. Text and blob results are bounded by a constant,
\texttt{Value::MAX\_BLOB\_LENGTH}, of $10^9$ bytes. A statement whose result
would exceed it should fail with the error \emph{string or blob too big}, which
is SQLite error code~18. Rust aborts the process when an allocation fails and cannot be
recovered, so a size check that runs after the allocation is not a check at all.

\paragraph{Differential testing.} A \emph{differential oracle} runs one query on
two engines and reports a disagreement. It needs no knowledge of the correct
answer, but it needs a \emph{reference engine} that implements the same feature.
Turso's fuzzer uses SQLite as its reference engine. SQLite has no \sql{lpad},
\sql{rpad} or \sql{repeat}, so a disagreement on those functions cannot be
judged.

\paragraph{Crash oracle.} A \emph{crash oracle} reports a statement that ends the
engine's process, panics, reports an internal error or does not answer within a
time limit. It does not report a \emph{user error}, which is a statement the
engine rejects with a normal message, such as a syntax error or a documented
limit. A crash oracle needs no reference engine. It finds no wrong answer that
leaves the process running.

\paragraph{Bug reports as test material.} DuckDB maintainers add the label
\emph{reproduced} to an issue once they have confirmed it. A fix usually adds a
regression test that contains the query which triggered the issue. Such a query
is small, it was chosen by someone who understood the defect, and it tends to sit
in an awkward corner: an empty string, a multi-byte character, a value at the
edge of a type.

\paragraph{Literal and template.} We call a constant in a query, such as
\sql{10} or \sql{'x'}, a \emph{literal}. A \emph{template} is a query in which
every literal is replaced by a typed hole, such as \sql{\{n0\}} for an integer or
\sql{\{s4\}} for a string. Filling the holes gives a new query with the same
structure.

\subsection{The Blind Spot of an Incumbent Fuzzer}
\label{sec:problem}

Let \Fz be the fuzzer that already tests an engine~$E$, and let \Rengine be the
reference engine that \Fz compares against. We write \Voc for the set of SQL
functions that \Fz can emit, and \Lit for the literals it can emit: a range of
integers and a set of strings. We write \VocR for the functions that \Rengine
implements.

A statement $q$ is \emph{expressible} by \Fz if every function in $q$ is in \Voc
and every literal in $q$ is in \Lit. A statement is \emph{judgeable} if every
function in $q$ is in \VocR. The \emph{blind spot} of \Fz is
\begin{equation}
  \label{eq:blind}
  \Blind = \{\, q \mid q \text{ is not expressible by } \Fz
           \ \lor\ q \text{ is not judgeable against } \Rengine \,\}.
\end{equation}
\Fz never generates a statement outside its own vocabulary or literal domain,
however long it runs. A statement that \Fz does generate but cannot judge is
generated in vain. Both halves of~\eqref{eq:blind} are properties of the source
code of \Fz and \Rengine, so \Blind can be measured before any test runs.

\paragraph{Goal.} Find statements $q \in \Blind$ on which $E$ violates the crash
oracle, without a reference engine, and with each finding regenerated from a
recorded seed. We look for two things: which functions and values put a
statement in \Blind, and where to get statements in that region that are likely
to break $E$.

\subsection{Why the Obvious Alternatives Fall Short}

\paragraph{Why not widen the incumbent fuzzer instead?} Its maintainers could add
\sql{lpad} to \Voc and \imax to \Lit, and the defects in this paper would then be
within reach. But that needs someone to know which function and which value to
add. The blind spot is large: any function outside \Voc and any value outside
\Lit belongs to it. Bug reports from another engine point at the corners where
someone has already been hurt. A crash oracle then removes the need for a
reference engine, so a function SQLite lacks can still be tested.

\paragraph{Why not run the other engine's tests directly?} A regression test was
written for the other engine's defect. Its literal values are tuned to that
defect, which the engine under test probably never had, and its syntax may not
parse. The structure of the query, meaning which functions are nested how, is
what carries over. The values have to be chosen again.
